\documentclass[11pt,compress,handout,blackandwhite]{beamer}

\usepackage{beamerthemesplit}
\usepackage[ansinew]{inputenc}
\usepackage[ngerman]{babel}
\usepackage{xcolor}


%%Pakete f�r mathematische Symbole und Umgebungen
\usepackage{amsmath,amsfonts,amssymb}

%%Paket f�r Bilder
\usepackage{epsfig}

\usepackage{marvosym}
%\usepackage{typearea}
%\usepackage{array}
%\usepackage{booktabs}
\usepackage{textcomp}

\definecolor{ubuntu}{rgb}{0.99,0.5,0.2}
\definecolor{dark-violet}{rgb}{0.40,0,0.48}
\definecolor{dunkelgrau.60}{gray}{0.40}




\title{Einf�hrung in die Shell}
\author{Pascal Bernhard}
\institute[BeLUG]{Berliner Linux User Group}
\date{\today}

\titlegraphic{\includegraphics[width=3.8cm,height=1.2cm]{belug-ev.png}}
\logo{\pgfimage[width=1.9cm,height=0.6cm]{Creative_Commons_BY-NC-SA.png}}

\begin{document}
%% Folie 1 (Titelfolie)
\frame{\titlepage}

\section[Vortrag]{}
\frame{\tableofcontents}

\section{Softwareverwaltung unter Linux}


%--------------------------------------------------------------------------------------%
%%%Folie 2

\subsection{1.Das Konzept Paketverwaltung}

\frame
{
  \frametitle{Was bedeutet Paketmanagement?}

  \begin{itemize}
  \item<1-> Linux-Software ist in Paketen organisiert
  \vspace{1.5cm}
  \item<2-> Pakete sind voneinander abh�ngig
	\begin{itemize}
	 \item \textcolor{dunkelgrau.60}{Programme teilen sich Bibliotheken}
	\end{itemize}
   
\vspace{1.5cm}

  \item<3-> Softwareinstallation erfolgt �ber Repositories

 

  \end{itemize}
}


%--------------------------------------------------------------------------------------%
%%%Folie 3


\subsection{2.Ubuntu Software-Center}

\frame
{
  \frametitle{Wie installiere ich Programme?}

  \begin{itemize}
  \item<1-> Ubuntu Software-Center

  \vspace{1.5cm}

  \item<2-> Kommandozeilentool: \textsl{\textcolor{dunkelgrau.60}{Apt}}

  \vspace{1.5cm}
  
  \item<3-> GUI f�r Apt: \textsl{\textcolor{dunkelgrau.60}{Synaptic}}
  \end{itemize}
}


%--------------------------------------------------------------------------------------%
%%%Folie 4


\subsection{3.Repositories}

\frame
{
  \frametitle{Was sind Repositories? Wo kommen die Pakete her?}
	
  \begin{itemize}
  \item<1-> Repositories sind Distributions-spezifisch

  \vspace{1.5cm}

  \item<2-> Konfiguration der Paketquellen in \texttt{\textcolor{dark-violet}{/etc/apt/sources.list}}

  \vspace{1.5cm}

  \item<3-> Hinzuf�gen von Repositories	
  \end{itemize}
}	
	
%--------------------------------------------------------------------------------------%
%%%Folie 5	
	
\subsection{4.Updates}

\frame
{
  \frametitle{Wie halte ich mein Ubuntu auf dem aktuellen Stand?}
  
  \begin{itemize}
  \item<2-> Updates/Upgrades werden zentral �ber das Paketmanagement gemacht f�r alle Pakete

  \vspace{1.5cm}

  \item<3-> Update �ber Synaptic

  \vspace{1.5cm}

  \item<4-> Update auf der Kommandozeile mit \texttt{\textcolor{ubuntu}{sudo apt-get update}} \& \texttt{\textcolor{ubuntu}{sudo apt-get upgrade}}
\end{itemize}
}

%--------------------------------------------------------------------------------------%
%%%Folie 6

\subsection{5.Installation/Entfernen von Paketen}

\begin{frame}[containsverbatim]
\frametitle{Wie installiere \& ich Pakete unter Ubuntu?}

  \begin{itemize}
   \item Update der Paketquellen \texttt{\textcolor{ubuntu}{sudo apt-get update}}

   \vspace{1.5cm}

   \item Suche nach Paketen \texttt{\textcolor{ubuntu}{sudo apt-cache search PAKETNAME}}
   
   \vspace{1.5cm}

   \item Installation von Paketen \texttt{\textcolor{ubuntu}{sudo apt-get install PAKETNAME1 PAKETNAME2}}
 
  \end{itemize}

\end{frame}

%--------------------------------------------------------------------------------------%
%%%Folie 7

\subsection{6.Sonderrepositories unter Ubuntu - PPAs}


\begin{frame}[containsverbatim]
\frametitle{PPAs \texttwelveudash Personal Package Archives}

  \begin{itemize}
   
    \item nicht-offizielle Repositories - keine Unterst�tzung durch Ubuntu/Cannonical

    \vspace{1.5cm}

    \item Beispiel: \textcolor{dunkelgrau.60}{Ubuntu Tweak}

      \begin{itemize}

       \item \texttt{\textcolor{ubuntu}{sudo add-apt-repository ppa:tualatrix/ppa}}

       \item \texttt{\textcolor{ubuntu}{sudo add-get update}}

       \item \texttt{\textcolor{ubuntu}{sudo apt-get install ubuntu-tweak}}

      \end{itemize}



  \end{itemize}



\end{frame}


%--------------------------------------------------------------------------------------%
%%%Folie 8

\subsection{7.Manuelle Installation von Paketen}







\end{document}
